\section{Discussão e limitações}
\label{sec:discussao}

Esta seção discute as limitações da reprodução em ordem de impacto. A
divergência de tamanho entre os CSVs públicos e o dataset do paper
(\S{}\ref{sec:discussao_dataset}) é a dominante; seguem a escolha de
não usar \texttt{clean\_df} (\S{}\ref{sec:discussao_clean}), a métrica
F1 ancorada no paper mas não publicada como chamada exata, a semente e
os hiperparâmetros não
publicados pelo paper, a contagem das ``22 \emph{features}'', a
substituição do SVM no \emph{combined}, a cobertura reduzida do
\emph{sweep} combinado e o alcance da generalização, encerrando com o
resumo das limitações.

\subsection{Divergência de tamanho é a limitação dominante}
\label{sec:discussao_dataset}

O \emph{achado} de \S{}\ref{sec:dataset_mismatch} domina todas as
demais limitações da reprodução: os CSVs públicos
\texttt{Train\_Test\_IoT\_*.csv} são um \emph{subset re-balanceado}
do dataset original. A classe \texttt{normal} foi sub-amostrada de
$\sim$35\,000 para $\sim$15\,000 amostras por dispositivo
(razão $\approx 0{,}43$), enquanto a classe \texttt{attack} foi
preservada na íntegra. Isso significa que mesmo o pipeline mais
fiel ao paper não consegue reproduzir os números exatos das
Tabelas 10--13, e que qualquer comparação entre nossa reprodução e o
paper deve ser lida como ``ordem de grandeza'', não como correspondência
exata.

A consequência concentra-se na tarefa \emph{combined} multi-classe.
Confrontando a Figura 5 do paper com os CSVs públicos, os \emph{ataques}
são idênticos linha a linha (\texttt{backdoor} 35\,k, \texttt{ddos}
25\,k, \texttt{injection} 35\,k, \texttt{password} 35\,k,
\texttt{ransomware} 16\,030, \texttt{scanning} 3\,973, \texttt{xss}
6\,116); a \emph{única} diferença é a classe \texttt{normal}
($245\,000$ no paper, $105\,000$ nos CSVs públicos). Como a acurácia e o
F1 \emph{weighted} multi-classe seguem a prevalência de \texttt{normal}
(a Tabela 13 do paper reporta acurácia de LR $=0{,}61$, igual à sua
fração de \texttt{normal}), nossos números ficam sistematicamente
abaixo: com $40\%$ de \texttt{normal} contra os $61\%$ do paper, há menos
da classe majoritária ``fácil'' para sustentar a métrica. A tarefa
\emph{combined} binária, ao contrário, reproduz bem
(\S{}\ref{sec:reproducao}): separar ataque de \texttt{normal} depende
pouco do volume da classe majoritária.

O \emph{Naïve Bayes} é o único modelo cujo \emph{gap} vai além da
prevalência. Nossa acurácia multi-classe ($0{,}106$) fica \emph{abaixo}
do \emph{baseline} de maioria ($0{,}40$), enquanto a do paper ($0{,}54$)
fica próxima da sua prevalência. A causa é uma patologia do
\texttt{GaussianNB}: as \emph{features} de proveniência (colunas de
dispositivo preenchidas com zero) e as variáveis quase-constantes por
classe produzem variâncias degeneradas que levam o modelo a prever
\texttt{ddos} em vez de \texttt{normal} (aumentar \texttt{var\_smoothing}
recupera apenas parcialmente, e foge do \texttt{GaussianNB} padrão do
paper). Acreditamos que os deltas residuais de \S{}\ref{sec:reproducao}
($|\Delta|<0{,}15$ em 91\% dos pares) são amplamente explicados pela
divergência de \texttt{normal}, e que sob o dataset original do paper a
aderência seria substancialmente maior.

\subsection{A escolha entre \emph{clean\_df} e não-\emph{clean\_df}}
\label{sec:discussao_clean}

O paper afirma, em sua \S{}VI-A1, que features categóricas são
encodadas como ``high$\rightarrow$1, low$\rightarrow$0''. Isso
sugere que o paper operou sobre \emph{features já limpas} (sem
variantes de \emph{whitespace}). Nos CSVs públicos, no entanto, as
quatro variantes (\texttt{'high'}, \texttt{'high\ '},
\texttt{'low'}, \texttt{'low\ '}) coexistem; o
\texttt{LabelEncoder} cru as trata como classes distintas.

Há \emph{duas} opções, e \textbf{nenhuma reproduz o paper
perfeitamente}:

\begin{itemize}[leftmargin=*,itemsep=0pt]
  \item \textbf{Com \texttt{clean\_df}}: as variantes são
        normalizadas; o sinal residual de encoding desaparece;
        F1 em \texttt{Fridge} e \texttt{Garage\_Door} cai para
        predição majoritária, e nossos números ficam muito
        \emph{abaixo} do paper.
  \item \textbf{Sem \texttt{clean\_df}} (escolhido): o encoding cru
        preserva qualquer sinal residual; o modelo aprende o que
        provavelmente aprendeu nos dados originais do paper;
        para algumas bases isso dá \emph{mais} sinal do que o paper
        teve (variantes \emph{whitespace} são informação extra),
        e nossos $\Delta$ ficam positivos.
\end{itemize}

Adotamos a segunda opção (a escolha que mais se aproxima de uma
reprodução honesta dada a versão de dataset que temos) e
documentamos explicitamente que isso pode inflar nossos números em
$\texttt{Fridge}$ para LR/SVM. O \texttt{clean\_df} permanece útil
em \S{}\ref{sec:dataset} para mostrar a magnitude do \emph{leak}
introduzido pelo encoding, mas não é parte do pipeline
\emph{paper-faithful}.

\subsection{Métrica F1: ancorada no paper, confirmada empiricamente}

A métrica adotada (\texttt{weighted}) está ancorada na definição de
média ponderada que o paper dá na discussão da Tabela 13
(\S{}\ref{sec:protocolo_metrica}); o que o paper não publica é a
\emph{chamada exata} de \texttt{sklearn} para as tarefas binárias
(Tabelas 10--11), que confirmamos empiricamente pela varredura das três
variantes. Se evidência adicional surgir (por exemplo, o código do paper
sendo publicado), poderemos re-rodar com a métrica exata. Nesse cenário, todos os
resultados desta versão podem ser recomputados em poucas horas
sobre os \emph{JSONL}s, sem re-treinar modelos: \texttt{f1} (binário),
\texttt{f1\_macro} e \texttt{f1\_weighted} são gravadas
simultaneamente em cada execução.

\subsection{Semente e hiperparâmetros não publicados pelo paper}

O paper não publica a semente de aleatoriedade do \emph{split} 80/20,
dos \emph{folds} nem da inicialização de RF/SVM/LSTM; fixamos
\texttt{random\_state=42} em todos os nossos experimentos, sem garantia
de paridade com a partição original. Além disso, o paper especifica
apenas os hiperparâmetros de cabeçalho de cada modelo (RF com 10
árvores e critério Gini, kNN com $k=5$, SVM RBF com
\texttt{gamma='auto'}) e deixa os demais (\texttt{C}, \texttt{solver},
\texttt{max\_iter}, \texttt{class\_weight}) nos \emph{defaults} do
\texttt{scikit-learn} 0.22.1 da época. Como esses \emph{defaults}
mudaram entre versões, o uso do \texttt{scikit-learn} atual é uma
fonte adicional e inevitável de divergência.

\subsection{Combined: as ``22 features'' do paper são contagem de colunas}

O paper \S{}VI-C reporta que o \emph{combined} foi montado ``\emph{with
a total of 22 features}''. Nosso \emph{combined} (via
\texttt{build\_combined\_dataset}) tem 17 \emph{features} de sensor,
o \texttt{UNION} das colunas das sete bases. A diferença é
\emph{contábil}, não conceitual: as 22 são a contagem de colunas do CSV
combinado, isto é, as 17 \emph{features} de sensor mais
\texttt{date}, \texttt{time}, \texttt{timestamp}, \texttt{label} e
\texttt{type}. O próprio paper, em sua \S{}VI-A1, declara descartar
\texttt{date}, \texttt{time} e \texttt{timestamp} do vetor de
\emph{features} (``\emph{were omitted\ldots\ as they may cause some ML
methods to overfit}''), exatamente como fazemos via
\texttt{drop\_temporal\_leak}. Logo, \textbf{ambos os pipelines treinam
sobre o mesmo conjunto de $\sim$17 \emph{features} de sensor}: não há
discrepância de \emph{features} de modelagem, e em particular o paper
\emph{não} obteve sinal extra por reter colunas temporais.

Em nossa construção do \emph{combined}, as numéricas faltantes do
\emph{union} das sete bases são preenchidas com zero (verificamos que
trocar a imputação por mediana não altera os F1 dos \emph{baselines}).
As \emph{features} categóricas, ao terem \texttt{NaN} convertido em um
\emph{label} próprio pelo \texttt{LabelEncoder}, viram um proxy de
``dispositivo de origem'': apenas \texttt{Fridge} preenche
\texttt{temp\_condition}, apenas \texttt{Garage\_Door} preenche
\texttt{sphone\_signal}, e assim por diante. \emph{Esse} \emph{leak} de
proveniência não estava no protocolo do paper e é uma fonte adicional
de divergência, especialmente para os modelos mais ávidos por
encoding categórico (LR, kNN, NB).

\subsection{Substituição do SVM RBF no \emph{combined}}

O SVM com kernel RBF é $\mathcal{O}(n^2)$ em tempo de treinamento,
inviável em $\sim$210\,000 linhas do \emph{combined}. Como o SVM é um
dos oito \emph{baselines} do paper e queremos preservá-lo na
comparação, substituímos apenas o \emph{kernel} RBF por
\texttt{LinearSVC} para $n>50$\,k, registrando o substituto na coluna
\texttt{effective\_model} do JSONL.

\subsection{Cobertura limitada do \emph{sweep} combined}

Os experimentos \emph{per-device} cobrem todo o produto cartesiano de
hiperparâmetros de cada modelo. Os experimentos \emph{combined}
usam apenas uma configuração canônica por modelo (geralmente a
melhor configuração descoberta no \emph{per-device}), por motivo de
custo computacional. Isso significa que os F1 \emph{combined}
reportados são limites inferiores: um \emph{sweep} completo poderia
encontrar configurações ligeiramente melhores.

\subsection{Generalização}

Todos os resultados acima são específicos aos CSVs públicos atuais
do TON\_IoT. Reprodutibilidade exata requer a versão original do
dataset (não publicamente disponível) ou que os autores publiquem o
\emph{seed} de re-amostragem que produziu os CSVs públicos.
Generalização para outros datasets IoT requer revalidar pelo menos
o pipeline de auditoria (Cramér's V, leak temporal) antes de
declarar qualquer ranking entre modelos.

\subsection{Resumo das limitações}

\begin{itemize}[leftmargin=*,itemsep=0pt]
  \item Versão de dataset diferente do paper (classe \texttt{normal}
        sub-amostrada para 43\% do volume original).
  \item Métrica F1 ancorada na média ponderada definida pelo paper
        (Tabela 13) e confirmada empiricamente; o paper não publica a
        chamada exata de \texttt{sklearn}.
  \item Semente de aleatoriedade não publicada pelo paper; fixamos
        \texttt{random\_state=42}, sem paridade garantida com a
        partição original.
  \item Hiperparâmetros não-cabeçalho dos \emph{baselines} deixados
        nos \emph{defaults} do \texttt{scikit-learn}, que diferem
        entre versões.
  \item Pipeline \emph{paper-faithful} sem \texttt{clean\_df}, que
        preserva \emph{leaks} do encoding (decisão honesta, mas
        documentada).
  \item Leak de proveniência no \emph{combined} (\texttt{NaN}
        categórico vira proxy de dispositivo de origem).
  \item SVM RBF substituído por \texttt{LinearSVC} no
        \emph{combined}.
  \item Cobertura reduzida do \emph{sweep} combined.
\end{itemize}
